% --- Archon physics formalization source begin ---
% archon:physics
% archon:covers PhyXMiniProblems/problem_phyx_mini_0577.lean
% archon:source-report reports/phyx_mini/problem_phyx_mini_0577.source.json

\chapter{Physics problem phyx\_mini\_0577}
\label{ch:PhyXMiniProblems_problem_phyx_mini_0577}

\paragraph{Problem source.}
\#\# Physical scenario

Comet stimulated emission. When a comet approaches the Sun, the increased warmth evaporates water from the ice on the surface of the comet nucleus, producing a thin atmosphere of water vapor around the nucleus. Sunlight can then dissociate H\textbackslash{}\_2O molecules in the vapor to H atoms and OH molecules. The sunlight can also excite the OH molecules to higher energy levels. When the comet is still relatively far from the Sun, the sunlight causes equal excitation to the \textbackslash{}( E\textbackslash{}\_2 \textbackslash{}) and \textbackslash{}( E\textbackslash{}\_1 \textbackslash{}) levels. Hence, there is no population inversion between the two levels. However, as the comet approaches the Sun, the excitation to the \textbackslash{}( E\textbackslash{}\_1 \textbackslash{}) level decreases and population inversion occurs. The reason has to do with Fraunhofer lines. As a comet approaches the Sun, the Doppler effect due to the comet's speed relative to the Sun shifts the Fraunhofer lines in wavelength, overlapping one of them with the wavelength required for excitation to the \textbackslash{}( E\textbackslash{}\_1 \textbackslash{}) level in OH molecules. Population inversion occurs in those molecules, and they radiate stimulated emission at about \textbackslash{}( 1666\textbackslash{}text\textbackslash{}\{ MHz \textbackslash{}\} \textbackslash{}).

\#\# Question

What was the energy difference \textbackslash{}( E\textbackslash{}\_2 - E\textbackslash{}\_1 \textbackslash{}) for that emission?

\#\# Answer choices

- A: A: 12.80μeV

- B: B: 1.67μeV

- C: C: 6.87μeV

- D: D: 3.43μeV

\#\# Figure caption (auxiliary; use the image as primary evidence)

The image displays a diagram illustrating energy levels and transitions between them. There are three horizontal red lines labeled as \textbackslash{}(E\_0\textbackslash{}), \textbackslash{}(E\_1\textbackslash{}), and \textbackslash{}(E\_2\textbackslash{}), representing different energy states. Two black arrows point upward, indicating transitions:

1. The first arrow points from \textbackslash{}(E\_0\textbackslash{}) to \textbackslash{}(E\_1\textbackslash{}).
2. The second arrow points from \textbackslash{}(E\_1\textbackslash{}) to \textbackslash{}(E\_2\textbackslash{}).

The figure represents energy absorption processes in a system, such as an atom or molecule, moving from a lower energy state to a higher one.

\#\# Dataset metadata

Quantum Phenomena; Physical Model Grounding Reasoning, Multi-Formula Reasoning

\paragraph{Recorded answer/context.}
C: C: 6.87μeV

\paragraph{Figure/image path.}
/root/proposal\_for\_physic/science-mango/phyx\_mini\_run/phyx\_data/test\_image/577.png

\paragraph{Formalization target.}
create a compiling Lean file with sorry bodies at `PhyXMiniProblems/problem\_phyx\_mini\_0577.lean`. The Lean declarations must preserve the physical quantities, dimensions or dimensional roles, figure labels, governing-law hypotheses, and final relation expressed by this problem.
Use Mathlib/Physlib names found through LeanExplore where available. If a physics API is missing, introduce faithful local abstractions rather than scalar placeholder aliases.

\begin{theorem}[Physics formalization target]
\label{thm:physics:phyx_mini_0577:target}
\lean{PhyXMiniProblems.ProblemPhyXMini0577.problem_phyx_mini_0577}
\uses{def:physics:phyx-mini-0577:phyxminiproblems-problemphyxmini0577-cometohemissionsetup, def:physics:phyx-mini-0577:phyxminiproblems-problemphyxmini0577-withinabsolutetolerance, def:physics:phyx-mini-0577:phyxminiproblems-problemphyxmini0577-matchescometohscenario, def:physics:phyx-mini-0577:phyxminiproblems-problemphyxmini0577-matchesprimaryenergylevelfigure, def:physics:phyx-mini-0577:phyxminiproblems-problemphyxmini0577-matchesfraunhoferdopplermechanism, def:physics:phyx-mini-0577:phyxminiproblems-problemphyxmini0577-hasphysicalohemissionparameters, def:physics:phyx-mini-0577:phyxminiproblems-problemphyxmini0577-usesstandardplanckconstant, def:physics:phyx-mini-0577:phyxminiproblems-problemphyxmini0577-satisfiesstimulatedemissionlaws, def:physics:phyx-mini-0577:phyxminiproblems-problemphyxmini0577-e2minuse1inmicroelectronvolts, def:physics:phyx-mini-0577:phyxminiproblems-problemphyxmini0577-recordeddatasetanswer, def:physics:phyx-mini-0577:phyxminiproblems-problemphyxmini0577-matchesdisplayedenergydifference, def:physics:phyx-mini-0577:phyxminiproblems-problemphyxmini0577-isuniqueclosestenergydifferencechoice}
The assigned autoformalize agent should translate this physics problem into Lean declarations in the covered file, with theorem and lemma proof bodies written as `by sorry`.
\end{theorem}
\begin{proof}
This is an autoformalization task, not a proof task. Produce statements that can later be proved by the physics prover without weakening the physical model.
\end{proof}

% --- Archon declaration topology begin ---
\section{Declaration topology}
\begin{definition}[energy Dimension]
  \label{def:physics:phyx-mini-0577:phyxminiproblems-problemphyxmini0577-energydimension}
  \lean{PhyXMiniProblems.ProblemPhyXMini0577.energyDimension}
  The MLT dimension of energy
\end{definition}
\begin{proof}
  This is the declared model definition of energy Dimension.
\end{proof}
\begin{definition}[action Dimension]
  \label{def:physics:phyx-mini-0577:phyxminiproblems-problemphyxmini0577-actiondimension}
  \lean{PhyXMiniProblems.ProblemPhyXMini0577.actionDimension}
  \uses{def:physics:phyx-mini-0577:phyxminiproblems-problemphyxmini0577-energydimension}
  Action has the dimension energy times time
\end{definition}
\begin{proof}
  This is the declared model definition of action Dimension.
\end{proof}
\begin{definition}[Frequency Quantity]
  \label{def:physics:phyx-mini-0577:phyxminiproblems-problemphyxmini0577-frequencyquantity}
  \lean{PhyXMiniProblems.ProblemPhyXMini0577.FrequencyQuantity}
  A nonnegative, unit-independent physical frequency
\end{definition}
\begin{proof}
  This is the declared model definition of Frequency Quantity.
\end{proof}
\begin{definition}[Wavelength Quantity]
  \label{def:physics:phyx-mini-0577:phyxminiproblems-problemphyxmini0577-wavelengthquantity}
  \lean{PhyXMiniProblems.ProblemPhyXMini0577.WavelengthQuantity}
  A nonnegative, unit-independent physical wavelength
\end{definition}
\begin{proof}
  This is the declared model definition of Wavelength Quantity.
\end{proof}
\begin{definition}[Planck Constant Quantity]
  \label{def:physics:phyx-mini-0577:phyxminiproblems-problemphyxmini0577-planckconstantquantity}
  \lean{PhyXMiniProblems.ProblemPhyXMini0577.PlanckConstantQuantity}
  \uses{def:physics:phyx-mini-0577:phyxminiproblems-problemphyxmini0577-actiondimension}
  A nonnegative physical value of the ordinary Planck constant h
\end{definition}
\begin{proof}
  This is the declared model definition of Planck Constant Quantity.
\end{proof}
\begin{definition}[energy In Joules]
  \label{def:physics:phyx-mini-0577:phyxminiproblems-problemphyxmini0577-energyinjoules}
  \lean{PhyXMiniProblems.ProblemPhyXMini0577.energyInJoules}
  Read a signed physical energy in coherent SI joules
\end{definition}
\begin{proof}
  This is the declared model definition of energy In Joules.
\end{proof}
\begin{definition}[frequency In Hertz]
  \label{def:physics:phyx-mini-0577:phyxminiproblems-problemphyxmini0577-frequencyinhertz}
  \lean{PhyXMiniProblems.ProblemPhyXMini0577.frequencyInHertz}
  \uses{def:physics:phyx-mini-0577:phyxminiproblems-problemphyxmini0577-frequencyquantity}
  Read a frequency in coherent SI hertz
\end{definition}
\begin{proof}
  This is the declared model definition of frequency In Hertz.
\end{proof}
\begin{definition}[frequency In Megahertz]
  \label{def:physics:phyx-mini-0577:phyxminiproblems-problemphyxmini0577-frequencyinmegahertz}
  \lean{PhyXMiniProblems.ProblemPhyXMini0577.frequencyInMegahertz}
  \uses{def:physics:phyx-mini-0577:phyxminiproblems-problemphyxmini0577-frequencyquantity, def:physics:phyx-mini-0577:phyxminiproblems-problemphyxmini0577-frequencyinhertz}
  Read a frequency in megahertz
\end{definition}
\begin{proof}
  This is the declared model definition of frequency In Megahertz.
\end{proof}
\begin{definition}[wavelength In Meters]
  \label{def:physics:phyx-mini-0577:phyxminiproblems-problemphyxmini0577-wavelengthinmeters}
  \lean{PhyXMiniProblems.ProblemPhyXMini0577.wavelengthInMeters}
  \uses{def:physics:phyx-mini-0577:phyxminiproblems-problemphyxmini0577-wavelengthquantity}
  Read a wavelength in coherent SI metres
\end{definition}
\begin{proof}
  This is the declared model definition of wavelength In Meters.
\end{proof}
\begin{definition}[speed In Meters Per Second]
  \label{def:physics:phyx-mini-0577:phyxminiproblems-problemphyxmini0577-speedinmeterspersecond}
  \lean{PhyXMiniProblems.ProblemPhyXMini0577.speedInMetersPerSecond}
  Read a speed in coherent SI metres per second
\end{definition}
\begin{proof}
  This is the declared model definition of speed In Meters Per Second.
\end{proof}
\begin{definition}[planck Constant In Joule Seconds]
  \label{def:physics:phyx-mini-0577:phyxminiproblems-problemphyxmini0577-planckconstantinjouleseconds}
  \lean{PhyXMiniProblems.ProblemPhyXMini0577.planckConstantInJouleSeconds}
  \uses{def:physics:phyx-mini-0577:phyxminiproblems-problemphyxmini0577-planckconstantquantity}
  Read a physical action in coherent SI joule-seconds
\end{definition}
\begin{proof}
  This is the declared model definition of planck Constant In Joule Seconds.
\end{proof}
\begin{definition}[Energy Level Label]
  \label{def:physics:phyx-mini-0577:phyxminiproblems-problemphyxmini0577-energylevellabel}
  \lean{PhyXMiniProblems.ProblemPhyXMini0577.EnergyLevelLabel}
  The three energy levels printed in the supplied diagram
\end{definition}
\begin{proof}
  This is the declared model definition of Energy Level Label.
\end{proof}
\begin{definition}[Chemical Species]
  \label{def:physics:phyx-mini-0577:phyxminiproblems-problemphyxmini0577-chemicalspecies}
  \lean{PhyXMiniProblems.ProblemPhyXMini0577.ChemicalSpecies}
  Chemical species explicitly named in the comet scenario
\end{definition}
\begin{proof}
  This is the declared model definition of Chemical Species.
\end{proof}
\begin{definition}[Comet Regime]
  \label{def:physics:phyx-mini-0577:phyxminiproblems-problemphyxmini0577-cometregime}
  \lean{PhyXMiniProblems.ProblemPhyXMini0577.CometRegime}
  The two comet regimes compared in the prose
\end{definition}
\begin{proof}
  This is the declared model definition of Comet Regime.
\end{proof}
\begin{definition}[Emission Mode]
  \label{def:physics:phyx-mini-0577:phyxminiproblems-problemphyxmini0577-emissionmode}
  \lean{PhyXMiniProblems.ProblemPhyXMini0577.EmissionMode}
  The emission mechanism named by the problem
\end{definition}
\begin{proof}
  This is the declared model definition of Emission Mode.
\end{proof}
\begin{definition}[Figure Arrow]
  \label{def:physics:phyx-mini-0577:phyxminiproblems-problemphyxmini0577-figurearrow}
  \lean{PhyXMiniProblems.ProblemPhyXMini0577.FigureArrow}
  The two black excitation arrows visible in the primary raster
\end{definition}
\begin{proof}
  This is the declared model definition of Figure Arrow.
\end{proof}
\begin{definition}[Comet Chemical Environment]
  \label{def:physics:phyx-mini-0577:phyxminiproblems-problemphyxmini0577-cometchemicalenvironment}
  \lean{PhyXMiniProblems.ProblemPhyXMini0577.CometChemicalEnvironment}
  \uses{def:physics:phyx-mini-0577:phyxminiproblems-problemphyxmini0577-chemicalspecies}
  Qualitative chemical content and sunlight interactions in the coma
\end{definition}
\begin{proof}
  This is the declared model definition of Comet Chemical Environment.
\end{proof}
\begin{definition}[Emission Transition]
  \label{def:physics:phyx-mini-0577:phyxminiproblems-problemphyxmini0577-emissiontransition}
  \lean{PhyXMiniProblems.ProblemPhyXMini0577.EmissionTransition}
  \uses{def:physics:phyx-mini-0577:phyxminiproblems-problemphyxmini0577-energylevellabel}
  A downward radiative transition from an initial to a final level
\end{definition}
\begin{proof}
  This is the declared model definition of Emission Transition.
\end{proof}
\begin{definition}[OHPrimary Energy Level Figure]
  \label{def:physics:phyx-mini-0577:phyxminiproblems-problemphyxmini0577-ohprimaryenergylevelfigure}
  \lean{PhyXMiniProblems.ProblemPhyXMini0577.OHPrimaryEnergyLevelFigure}
  \uses{def:physics:phyx-mini-0577:phyxminiproblems-problemphyxmini0577-energylevellabel, def:physics:phyx-mini-0577:phyxminiproblems-problemphyxmini0577-figurearrow}
  Literal line, label, and arrow information carried by the primary image
\end{definition}
\begin{proof}
  This is the declared model definition of OHPrimary Energy Level Figure.
\end{proof}
\begin{definition}[Comet OHEmission Setup]
  \label{def:physics:phyx-mini-0577:phyxminiproblems-problemphyxmini0577-cometohemissionsetup}
  \lean{PhyXMiniProblems.ProblemPhyXMini0577.CometOHEmissionSetup}
  \uses{def:physics:phyx-mini-0577:phyxminiproblems-problemphyxmini0577-frequencyquantity, def:physics:phyx-mini-0577:phyxminiproblems-problemphyxmini0577-wavelengthquantity, def:physics:phyx-mini-0577:phyxminiproblems-problemphyxmini0577-planckconstantquantity, def:physics:phyx-mini-0577:phyxminiproblems-problemphyxmini0577-energylevellabel, def:physics:phyx-mini-0577:phyxminiproblems-problemphyxmini0577-chemicalspecies, def:physics:phyx-mini-0577:phyxminiproblems-problemphyxmini0577-cometregime, def:physics:phyx-mini-0577:phyxminiproblems-problemphyxmini0577-emissionmode, def:physics:phyx-mini-0577:phyxminiproblems-problemphyxmini0577-cometchemicalenvironment, def:physics:phyx-mini-0577:phyxminiproblems-problemphyxmini0577-emissiontransition, def:physics:phyx-mini-0577:phyxminiproblems-problemphyxmini0577-ohprimaryenergylevelfigure}
  The independent physical state used by the model. In particular, the three energies and the emitted photon energy are fields; none is defined from an answer choice or from the requested 6.87 micro-eV value
\end{definition}
\begin{proof}
  This is the declared model definition of Comet OHEmission Setup.
\end{proof}
\begin{definition}[Within Absolute Tolerance]
  \label{def:physics:phyx-mini-0577:phyxminiproblems-problemphyxmini0577-withinabsolutetolerance}
  \lean{PhyXMiniProblems.ProblemPhyXMini0577.WithinAbsoluteTolerance}
  Absolute-error agreement of a scalar readout with a nominal value
\end{definition}
\begin{proof}
  This is the declared model definition of Within Absolute Tolerance.
\end{proof}
\begin{definition}[Matches Comet OHScenario]
  \label{def:physics:phyx-mini-0577:phyxminiproblems-problemphyxmini0577-matchescometohscenario}
  \lean{PhyXMiniProblems.ProblemPhyXMini0577.MatchesCometOHScenario}
  \uses{def:physics:phyx-mini-0577:phyxminiproblems-problemphyxmini0577-frequencyinmegahertz, def:physics:phyx-mini-0577:phyxminiproblems-problemphyxmini0577-cometohemissionsetup, def:physics:phyx-mini-0577:phyxminiproblems-problemphyxmini0577-withinabsolutetolerance}
  Direct transcription of the physical narrative. The one-megahertz window makes the word "about" explicit; it is an observational premise, not an assumption about the energy gap sought in the question
\end{definition}
\begin{proof}
  This is the declared model definition of Matches Comet OHScenario.
\end{proof}
\begin{definition}[Matches Primary Energy Level Figure]
  \label{def:physics:phyx-mini-0577:phyxminiproblems-problemphyxmini0577-matchesprimaryenergylevelfigure}
  \lean{PhyXMiniProblems.ProblemPhyXMini0577.MatchesPrimaryEnergyLevelFigure}
  \uses{def:physics:phyx-mini-0577:phyxminiproblems-problemphyxmini0577-ohprimaryenergylevelfigure}
  The primary bitmap has three horizontal red level lines and two upward black arrows. Contrary to the auxiliary caption, both arrows start at E0
\end{definition}
\begin{proof}
  This is the declared model definition of Matches Primary Energy Level Figure.
\end{proof}
\begin{definition}[Matches Fraunhofer Doppler Mechanism]
  \label{def:physics:phyx-mini-0577:phyxminiproblems-problemphyxmini0577-matchesfraunhoferdopplermechanism}
  \lean{PhyXMiniProblems.ProblemPhyXMini0577.MatchesFraunhoferDopplerMechanism}
  \uses{def:physics:phyx-mini-0577:phyxminiproblems-problemphyxmini0577-wavelengthinmeters, def:physics:phyx-mini-0577:phyxminiproblems-problemphyxmini0577-speedinmeterspersecond, def:physics:phyx-mini-0577:phyxminiproblems-problemphyxmini0577-cometohemissionsetup}
  Qualitative Doppler/Fraunhofer mechanism stated in the problem: approach produces a shorter shifted wavelength and that line overlaps the wavelength needed for excitation to E1. No numerical energy answer occurs here
\end{definition}
\begin{proof}
  This is the declared model definition of Matches Fraunhofer Doppler Mechanism.
\end{proof}
\begin{definition}[Has Physical OHEmission Parameters]
  \label{def:physics:phyx-mini-0577:phyxminiproblems-problemphyxmini0577-hasphysicalohemissionparameters}
  \lean{PhyXMiniProblems.ProblemPhyXMini0577.HasPhysicalOHEmissionParameters}
  \uses{def:physics:phyx-mini-0577:phyxminiproblems-problemphyxmini0577-energyinjoules, def:physics:phyx-mini-0577:phyxminiproblems-problemphyxmini0577-frequencyinhertz, def:physics:phyx-mini-0577:phyxminiproblems-problemphyxmini0577-wavelengthinmeters, def:physics:phyx-mini-0577:phyxminiproblems-problemphyxmini0577-planckconstantinjouleseconds, def:physics:phyx-mini-0577:phyxminiproblems-problemphyxmini0577-cometohemissionsetup}
  Positivity and ordering conditions selecting a physical OH-level model
\end{definition}
\begin{proof}
  This is the declared model definition of Has Physical OHEmission Parameters.
\end{proof}
\begin{definition}[Uses Standard Planck Constant]
  \label{def:physics:phyx-mini-0577:phyxminiproblems-problemphyxmini0577-usesstandardplanckconstant}
  \lean{PhyXMiniProblems.ProblemPhyXMini0577.UsesStandardPlanckConstant}
  \uses{def:physics:phyx-mini-0577:phyxminiproblems-problemphyxmini0577-planckconstantinjouleseconds, def:physics:phyx-mini-0577:phyxminiproblems-problemphyxmini0577-cometohemissionsetup}
  Exact SI calibration of the independently known ordinary Planck constant
\end{definition}
\begin{proof}
  This is the declared model definition of Uses Standard Planck Constant.
\end{proof}
\begin{definition}[Satisfies Stimulated Emission Laws]
  \label{def:physics:phyx-mini-0577:phyxminiproblems-problemphyxmini0577-satisfiesstimulatedemissionlaws}
  \lean{PhyXMiniProblems.ProblemPhyXMini0577.SatisfiesStimulatedEmissionLaws}
  \uses{def:physics:phyx-mini-0577:phyxminiproblems-problemphyxmini0577-energyinjoules, def:physics:phyx-mini-0577:phyxminiproblems-problemphyxmini0577-frequencyinhertz, def:physics:phyx-mini-0577:phyxminiproblems-problemphyxmini0577-planckconstantinjouleseconds, def:physics:phyx-mini-0577:phyxminiproblems-problemphyxmini0577-cometohemissionsetup}
  The two governing stimulated-emission relations. They apply to the selected transition but contain no numerical energy gap or answer label
\end{definition}
\begin{proof}
  This is the declared model definition of Satisfies Stimulated Emission Laws.
\end{proof}
\begin{definition}[e2 Minus E1 In Micro Electron Volts]
  \label{def:physics:phyx-mini-0577:phyxminiproblems-problemphyxmini0577-e2minuse1inmicroelectronvolts}
  \lean{PhyXMiniProblems.ProblemPhyXMini0577.e2MinusE1InMicroElectronVolts}
  \uses{def:physics:phyx-mini-0577:phyxminiproblems-problemphyxmini0577-energyinjoules, def:physics:phyx-mini-0577:phyxminiproblems-problemphyxmini0577-cometohemissionsetup}
  Energy-level gap E2 - E1, read in micro-electron-volts
\end{definition}
\begin{proof}
  This is the declared model definition of e2 Minus E1 In Micro Electron Volts.
\end{proof}
\begin{definition}[Answer Choice]
  \label{def:physics:phyx-mini-0577:phyxminiproblems-problemphyxmini0577-answerchoice}
  \lean{PhyXMiniProblems.ProblemPhyXMini0577.AnswerChoice}
  Labels of the four choices printed with the problem
\end{definition}
\begin{proof}
  This is the declared model definition of Answer Choice.
\end{proof}
\begin{definition}[displayed Energy Difference In Micro Electron Volts]
  \label{def:physics:phyx-mini-0577:phyxminiproblems-problemphyxmini0577-displayedenergydifferenceinmicroelectronvolts}
  \lean{PhyXMiniProblems.ProblemPhyXMini0577.displayedEnergyDifferenceInMicroElectronVolts}
  \uses{def:physics:phyx-mini-0577:phyxminiproblems-problemphyxmini0577-answerchoice}
  Micro-electron-volt value printed beside each answer label
\end{definition}
\begin{proof}
  This is the declared model definition of displayed Energy Difference In Micro Electron Volts.
\end{proof}
\begin{definition}[recorded Dataset Answer]
  \label{def:physics:phyx-mini-0577:phyxminiproblems-problemphyxmini0577-recordeddatasetanswer}
  \lean{PhyXMiniProblems.ProblemPhyXMini0577.recordedDatasetAnswer}
  \uses{def:physics:phyx-mini-0577:phyxminiproblems-problemphyxmini0577-answerchoice}
  Answer label recorded in the source dataset
\end{definition}
\begin{proof}
  This is the declared model definition of recorded Dataset Answer.
\end{proof}
\begin{definition}[Matches Displayed Energy Difference]
  \label{def:physics:phyx-mini-0577:phyxminiproblems-problemphyxmini0577-matchesdisplayedenergydifference}
  \lean{PhyXMiniProblems.ProblemPhyXMini0577.MatchesDisplayedEnergyDifference}
  \uses{def:physics:phyx-mini-0577:phyxminiproblems-problemphyxmini0577-cometohemissionsetup, def:physics:phyx-mini-0577:phyxminiproblems-problemphyxmini0577-withinabsolutetolerance, def:physics:phyx-mini-0577:phyxminiproblems-problemphyxmini0577-e2minuse1inmicroelectronvolts, def:physics:phyx-mini-0577:phyxminiproblems-problemphyxmini0577-answerchoice, def:physics:phyx-mini-0577:phyxminiproblems-problemphyxmini0577-displayedenergydifferenceinmicroelectronvolts}
  The 0.03 micro-eV tolerance accommodates the word "about" in 1666 MHz and the small mismatch between that rounded frequency and the displayed answer. It is far smaller than the separation between answer choices
\end{definition}
\begin{proof}
  This is the declared model definition of Matches Displayed Energy Difference.
\end{proof}
\begin{definition}[Is Unique Closest Energy Difference Choice]
  \label{def:physics:phyx-mini-0577:phyxminiproblems-problemphyxmini0577-isuniqueclosestenergydifferencechoice}
  \lean{PhyXMiniProblems.ProblemPhyXMini0577.IsUniqueClosestEnergyDifferenceChoice}
  \uses{def:physics:phyx-mini-0577:phyxminiproblems-problemphyxmini0577-cometohemissionsetup, def:physics:phyx-mini-0577:phyxminiproblems-problemphyxmini0577-e2minuse1inmicroelectronvolts, def:physics:phyx-mini-0577:phyxminiproblems-problemphyxmini0577-answerchoice, def:physics:phyx-mini-0577:phyxminiproblems-problemphyxmini0577-displayedenergydifferenceinmicroelectronvolts}
  The selected choice is strictly closer to the modeled gap than each rival
\end{definition}
\begin{proof}
  This is the declared model definition of Is Unique Closest Energy Difference Choice.
\end{proof}
% --- Archon declaration topology end ---
% --- Archon physics formalization source end ---
